\subsection{收敛级数}

3.1.1. 设 \(f = \sum f_{\alpha}\) 是属于 \(\hat{\mathbf{P}}(\mathbf{E}_1, \ldots, \mathbf{E}_n; \mathbf{F})\) 的形式级数（附录，第 A.5 号）。若 \(\gamma\) 是 \(\mathbf{F}\) 上的连续半范数，且 \(\mathbf{R} = (\mathbf{R}_1, \ldots, \mathbf{R}_n)\) 是 \(n\) 个严格正实数组成的序列，置：

\[
\| f \| _ {\gamma , \mathbf {R}} = \sum_ {\alpha} \mathbf {R} ^ {\alpha} \| f _ {\alpha} \| _ {\gamma}.
\]

若 \(F\) 是赋范空间且 \(\gamma\) 是 \(F\) 的范数，则用 \(\| f \|_{\mathbb{R}}\) 代替 \(\| f \|_{\gamma, \mathbb{R}}\)。

由满足对 \(\mathcal{H}_{\mathbb{R}}(\mathbf{E}_{1},\ldots,\mathbf{E}_{n};\mathbf{F})\) 上每个连续半范数 \(f\in\hat{\mathbf{P}}(\mathbf{E}_{1},\ldots,\mathbf{E}_{n};\mathbf{F})\)，\(\|f\|_{\gamma,\mathbb{R}}\) 都有限的 \(\gamma\) 构成的集合 \(\mathbf{F}\)，是 \(\hat{\mathbf{P}}(\mathbf{E}_{1},\ldots,\mathbf{E}_{n};\mathbf{F})\) 的向量子空间。有

\[
\mathcal {H} _ {\mathbb {R}} (\mathrm{E} _ {1}, \dots , \mathrm{E} _ {n}; \mathrm{F}) \subset \mathcal {H} _ {\mathbb {R} ^ {\prime}} (\mathrm{E} _ {1}, \dots , \mathrm{E} _ {n}; \mathrm{F})
\]

只要 \(R_{i} \geqslant R_{i}'\)（\(1 \leqslant i \leqslant n\)）即可。下列各空间的并

\[
\mathcal {H} _ {\mathbf {R}} (\mathrm{E} _ {1}, \dots , \mathrm{E} _ {n}; \mathrm{F})
\]

\(\mathcal{H}(E_{1},\ldots,E_{n};F)\)是 \(\hat{\mathrm{P}}(E_{1},\ldots,E_{n};F)\) 的向量子空间，记作 \(E_{i}\)，其中元素称为取值于 \(F\) 的各 \(E_{i}\) 的乘积上的收敛级数。它只取决于各 \(\mathcal{H}(E_{1},\ldots,E_{n};F)\) 的拓扑而不取决于它们的范数。因此，当各 \(E_{i}\) 是可赋范空间时，无须在每个 \(E_{i}\) 上选取范数，也可以谈论空间 H(E＿1,,E＿n;F)。

映射 \(f\mapsto\|f\|_{\gamma,\mathbb{R}}\) 是 \(\mathcal{H}_{\mathbb{R}}(\mathrm{E}_{1},\ldots,\mathrm{E}_{n};\mathrm{F})\) 上的半范数。当 \(\gamma\) 遍历 F 上的连续半范数（或仅遍历定义 F 的拓扑的一组半范数）时，由这些半范数定义的拓扑是豪斯多夫的。若 F 可赋范（分别地，完备），则 \(\mathcal{H}_{\mathbb{R}}(\mathrm{E}_{1},\ldots,\mathrm{E}_{n};\mathrm{F})\) 也可赋范（分别地，完备）。\(\mathcal{H}(\mathrm{E}_{1},\ldots,\mathrm{E}_{n};\mathrm{F})\) 上那些限制到每个 \(\mathcal{H}_{\mathbb{R}}\) 都连续的半范数，在 \(\mathcal{H}(\mathrm{E}_{1},\ldots,\mathrm{E}_{n};\mathrm{F})\) 上定义一个豪斯多夫拓扑，该拓扑只取决于各 \(\mathrm{E}_{i}\) 的拓扑。从 \(\mathcal{H}(\mathrm{E}_{1},\ldots,\mathrm{E}_{n};\mathrm{F})\) 到 \(\hat{\mathrm{P}}(\mathrm{E}_{1},\ldots,\mathrm{E}_{n};\mathrm{F})\) 的嵌入是连续的。

3.1.3. \(\hat{\mathbf{P}}(\mathbf{E};\mathbf{F})\) 到 \(\hat{\mathbf{P}}(\mathbf{E}_{1},\ldots,\mathbf{E}_{n};\mathbf{F})\) 的规范同构通过限制给出从 \(\mathcal{H}(\mathbf{E};\mathbf{F})\) 到 \(\mathcal{H}(\mathbf{E}_{1},\ldots,\mathbf{E}_{n};\mathbf{F})\) 的拓扑向量空间同构。

3.1.4. 设 \(f \in \mathcal{H}(\mathrm{E}_1, \ldots, \mathrm{E}_n; \mathrm{F})\)，并令 \(J(f)\)\(R \in (\mathbf{R}_+^*)^n\) 为满足 \(f \in \mathcal{H}_{\mathbb{R}}(\mathrm{E}_1, \ldots, \mathrm{E}_n; \mathrm{F})\) 的 \(I(f)\) 的集合。\(J(f)\) 的内部 I(\(f\)) 称为 f 的严格收敛指标。它由满足存在 \(R\)\(R' \in J(f)\)'  J(f) 且 \(0 < R_i < R'_i\)（\(1 \leqslant i \leqslant n\)）的 R 构成。用 \(\Omega(f)\)\((\log R_1, \ldots, \log R_n)\) 表示 R 属于 I(f) 时 \(R^n\) 中的点 \(R \in I(f)\) 的集合：它是 \(R^n\) 的凸子集。

当 \(n=1\) 时，集合 \(\mathbf{I}(f)\) 是 \([0,\rho(f)]\) 中的区间 \(\mathbf{R}\)，\(\rho(f)\) 称为 \(f\)\(+\infty\) 的严格收敛半径。它也是满足如下条件的实数 \(R>0\) 的集合的上确界（有限或 \(\gamma\)）：对 \(F\) 上每个连续半范数 ，存在常数 \(M\)，使得（其中 \(\|f_{m}\|_{\gamma}\leqslant MR^{-m}\)）对每个整数 \(f=\sum f_{m},f_{m}\in P_{m}(E;F)\) 都有 \(m\geqslant0\)。

满足存在 \(x = (x_i) \in \mathbf{E}_1 \times \cdots \times \mathbf{E}_n\) 且对所有 i 有 \(\mathbf{R} \in \mathbf{I}(f)\) 的点 \(\|x_i\| \leqslant R_i\) 的集合称为 f 的严格收敛域，记作 \(\mathbf{C}(f)\)。它也是满足存在 \(\mathbf{R} \in \mathbf{J}(f)\) 且对所有 i 有 \(\|x_i\| < R_i\) 的点 x 的集合的内部。

3.1.5. 对于 \(f_{\alpha} \in \mathbf{P}_{\alpha}(\mathbf{E}_1, \ldots, \mathbf{E}_n; \mathbf{F})\) 以及 \(\gamma\) 上每个连续半范数 \(\mathbf{F}\)，置：

\[
\| f _ {\alpha} \| _ {\gamma} = \sup _ {\| x _ {i} \| \leqslant 1} \| f _ {\alpha} (x _ {1}, \dots , x _ {n}) \| _ {\gamma}.
\]

于是有不等式：

\[
\left\| \left| f _ {\alpha} \right| \right\| _ {\gamma} \leqslant \left\| f _ {\alpha} \right\| _ {\gamma} \leqslant \frac {\alpha^ {\alpha}}{\alpha !} \left\| \left| f _ {\alpha} \right| \right\| _ {\gamma}.
\]

对于 \(f = \sum f_{\alpha} \in \hat{\mathbf{P}}(\mathrm{E}_1, \ldots, \mathrm{E}_n; \mathbf{F})\) 和 \(\mathbf{R} \in (\mathbf{R}_{+}^{*})^n\)，置：

\[
\| | f \| | _ {\gamma , \mathbf {R}} = \sum_ {\alpha} \mathbf {R} ^ {\alpha} \| | f _ {\alpha} \| | _ {\gamma}
\]

并用 \(\tilde{\mathcal{H}}_{\mathbb{R}}(\mathrm{E}_{1},\ldots ,\mathrm{E}_{n};\mathrm{F})\) 表示下列空间的向量子空间

\[
\hat {\mathbf {P}} (\mathbf {E} _ {1}, \dots , \mathbf {E} _ {n}; \mathbf {F})
\]

由满足对每个  都有  | \(f\)\(\| | f \| | _ {\gamma,\mathbb{R}}\)  | ＿ \(\gamma\),R 有限的 f 构成，并赋予由半范数 \(f\mapsto\| | f \| | _ {\gamma,\mathbb{R}}\) 定义的拓扑。有 \(\mathcal{H}_{\mathbb{R}}\subset\tilde{\mathcal{H}}_{\mathbb{R}}\)，且从 \(\mathcal{H}_{\mathbb{R}}\) 到 \(\tilde{\mathcal{H}}_{\mathbb{R}}\) 的嵌入是连续的。空间 \(\mathcal{H}(\mathrm{E}_{1},\ldots,\mathrm{E}_{n};\mathrm{F})\) 是各 \(\tilde{\mathcal{H}}_{\mathbb{R}}(\mathrm{E}_{1},\ldots,\mathrm{E}_{n};\mathrm{F})\) 的并，其拓扑是使 \(\mathcal{H}_{\mathbb{R}}\) 到 \(\mathcal{H}\) 的嵌入连续的最细局部凸拓扑。

若 \(f \in \mathcal{H}(E_1, \ldots, E_n; F)\)，则称满足 \(\tilde{I}(f)\) 的 \(\tilde{J}(f)\) 的集合 \(R \in (\mathbf{R}_+^*)\) 的内部 \(f \in \mathcal{H}_{\mathbb{R}}\) 为收敛指标。有：

\[
e ^ {- 1} \mathbf {I} (f) \subset \tilde{\mathbf {I}} (f) \subset \mathbf {I} (f)
\]

（其中 e 是自然对数的底）。从 \(\tilde{I}(f)\) 出发，按 3.1.4 定义收敛域 \(\tilde{C}(f)\)，并在 \(n = 1\) 时定义收敛半径 \(\tilde{\rho}(f)\)。特别地，有：

\[
e ^ {- 1} \tilde {\rho} (f) \leqslant \rho (f) \leqslant \tilde {\rho} (f).
\]

3.1.6. 对于 \(\mathbf{R} \in (\mathbf{R}_{+}^{*})^{n}\)，将 E 中以 0 为中心、半径为 \(\mathbf{R}\) 的（闭）多球称为由满足对所有 i 都有 \(\mathbf{B}(\mathbf{R})\) 的 \(x \in \mathbf{E}\) 构成的集合 \(\|x_{i}\| \leqslant \mathbf{R}_{i}\)。若  E＿\(i\)\(\dim \mathbf{E}_{i} = 1\) = 1，也称为多圆盘。若 \(f \in \mathcal{H}(\mathbf{E}_{1}, \ldots, \mathbf{E}_{n}; \mathbf{F})\)，则 \(f\) 的收敛域（分别地，严格收敛域）是 \(\mathbf{B}(\mathbf{R})\)（分别地，\(\mathbf{R} \in \tilde{\mathbf{I}}(f)\)）时多球 \(\mathbf{R} \in \mathbf{I}(f)\) 的并。

3.1.7. 设 \(f \in \mathcal{H}(E_1, \ldots, E_n; F)\)。假设 F 拟完备。对每个 \(x \in \tilde{\mathbf{C}}(f)\)，族 \(f_{\alpha}(x)\) 在 F 中可求和。其和记作 \(\hat{f}(x)\)，或简写为 \(f(x)\)，是 \(\tilde{\mathbf{C}}(f)\) 上的连续函数。更确切地说，对每个满足 \(f \in \tilde{\mathcal{H}}_{\mathbb{R}}\) 的 R，族 \(f_{\alpha}(x)\) 当 \(x \in B(R)\) 时一致可求和。映射 \(f \mapsto \hat{f}\) 是从 \(\tilde{\mathcal{H}}_{\mathbb{R}}\) 到 \(B(R)\) 上有界连续函数空间的连续线性单射，后者赋予一致收敛拓扑。

3.1.8. 设 \(F_1, \ldots, F_m\) 是分离多范数空间，\(u\)\(m\) 是从 \(F_1 \times \cdots \times F_m\) 到 \(F\) 的连续 m 线性映射。设 \(f_i \in \mathcal{H}(E_1, \ldots, E_n; F_i)\)（\(1 \leqslant i \leqslant m\)）。形式级数 \(u(f_1, \ldots, f_m)\) 属于 \(\mathcal{H}(E_1, \ldots, E_n; F)\)。有：

\[
\mathrm{C} (u (f _ {1}, \dots , f _ {m})) \supset \bigcap_ {i} \mathrm{C} (f _ {i})
\]

\[
\tilde {C} (u (f _ {1}, \dots , f _ {m})) \supset \bigcap_ {i} \tilde {C} (f _ {i})
\]

并且，若 \(x \in \bigcap_{i} \tilde{C}(f_i)\)，且 F 和各 \(F_i\) 拟完备，则有：

\[
u (f _ {1}, \dots , f _ {m}) (x) = u (f _ {1} (x), \dots , f _ {m} (x)).
\]

3.1.9. 设 \(\mathbf{F}_1, \ldots, \mathbf{F}_m\) 是完备赋范空间，且 \(\mathbf{F}\) 拟完备。令 \(\mathbf{f} = (f_i)_{1 \leqslant i \leqslant m}\)，其中 \(f_i \in \mathcal{H}(\mathbf{E}_1, \ldots, \mathbf{E}_n; \mathbf{F}_i)\)，\(g \in \mathcal{H}(\mathbf{F}_1, \ldots, \mathbf{F}_m; \mathbf{F})\)，并假设 \((f_i(0))_{1 \leqslant i \leqslant m}\) 属于 \(g\) 的严格收敛域。于是，对每个 \(\alpha \in \mathbb{N}^m\)，形式级数 \(g_\alpha \circ \mathbf{f}\) 属于 \(\mathcal{H}(\mathbf{E}_1, \ldots, \mathbf{E}_n; \mathbf{F})\)，且族 \(g_\alpha \circ \mathbf{f}\) 在 \(\mathcal{H}(\mathbf{E}_1, \ldots, \mathbf{E}_n; \mathbf{F})\) 中可求和，从而当然也在 \(\hat{\mathbf{P}}(\mathbf{E}_1, \ldots, \mathbf{E}_n; \mathbf{F})\) 中可求和。其和记作 \(g \circ \mathbf{f}\)。

更确切地说，存在 \(\mathbb{R}\in\bigcap_{i}\mathrm{I}(f_{i})\) 和 \(\mathbb{R}^{\prime}\in\mathrm{I}(g)\)，使得当 \(\|f_{i}\|_{\mathbb{R}}<R_{i}^{\prime}\) 时 \(1\leqslant i\leqslant m\)。在这些条件下，形式级数 \(g_{\alpha}\circ f\) 属于 \(\mathcal{H}_{\mathbb{R}}(\mathrm{E}_{1},\ldots,\mathrm{E}_{n};\mathrm{F})\)，且族 \(g_{\alpha}\circ f\) 在 \(\mathcal{H}_{\mathbb{R}}(\mathrm{E}_{1},\ldots,\mathrm{E}_{n};\mathrm{F})\) 中可求和。最后，若 \(x\in\mathrm{B}(\mathbb{R})\)，则 \(\mathbf{f}(x)=(f_{i}(x))\) 属于 \(\mathrm{B}(\mathrm{R}^{\prime})\subset\mathrm{F}_{1}\times\cdots\times\mathrm{F}_{m}\)，且有：

\[
g (\mathbf {f} (x)) = (g \circ \mathbf {f}) (x).
\]

3.1.10. 假设 \(E_{i} = K\)（\(1 \leqslant i \leqslant n\)）。于是将空间 \(\hat{\mathbf{P}}(\mathbf{K}^{n}; \mathbf{F})\) 识别为具有 F 中系数的 n 个不定元 \(X_{1}, \ldots, X_{n}\) 的形式幂级数空间，且 \(\hat{\mathbf{P}}(\mathbf{K}^{n};\mathbf{F})\) 中的元素写作：

\[
f = \sum_ {\alpha} X ^ {\alpha} c _ {\alpha} \quad \text {   with   } c _ {\alpha} \in F.
\]

若 \(\mathbf{R} \in (\mathbf{R}_{+}^{*})^{n}\)，且 \(\gamma\) 是 \(\mathbf{F}\) 上的连续半范数，则有：

\[
\| | f \| | _ {\gamma , \mathbf {R}} = \| f \| _ {\gamma , \mathbf {R}} = \sum_ {\alpha} \mathbf {R} ^ {\alpha} \| c _ {\alpha} \| _ {\gamma}
\]

有 \(\mathbf{I}(f) = \tilde{\mathbf{I}}(f)\)、\(\mathbf{C}(f) = \tilde{\mathbf{C}}(f)\)，并且当 \(n = 1\) 时 \(\rho(f) = \tilde{\rho}(f)\)。

系数在 \(\mathcal{H}(\mathbf{K}^{n};\mathbf{K})\) 中的收敛级数空间 \(\mathbf{K}\) 也记作 \(\mathbf{K}\{\{\mathbf{X}_{1},\ldots ,\mathbf{X}_{n}\}\}\)；它是 \(\mathbf{K}[[\mathbf{X}_1,\dots ,\mathbf{X}_n]]\) 的子代数。空间 \(\mathcal{H}(\mathbf{K}^n;\mathbf{F})\) 是 \(\mathbf{K}\{\{\mathbf{X}_1,\ldots ,\mathbf{X}_n\}\}\) 上的模；若 \(\mathbf{F}\) 有限维，则将此模识别为 \(\mathbf{K}\{\{\mathbf{X}_1,\ldots ,\mathbf{X}_n\}\} \otimes_{\mathbf{K}}\mathbf{F}\)。

3.1.11. 设 \(\mathbf{f} \in \mathcal{H}(\mathbf{K}^n; \mathbf{K}^m)\)，它由系数在 K 中的 \(m\) 个收敛级数 \(f_j(X_1, \ldots, X_n)\) 组成的系统表示。同样，设 \(\mathbf{g} \in \mathcal{H}(\mathbf{K}^m; \mathbf{K}^p)\)，它由 \(p\) 个收敛级数 \(g_k(Y_1, \ldots, Y_m) = \sum g_{k,\beta} Y^\beta\) 组成的系统表示。空间 \(\mathbf{h} = \mathbf{g} \circ \mathbf{f}\) 中的元素 \(\mathcal{H}(\mathbf{K}^n; \mathbf{K}^p)\)（参见 3.1.9）由 \(p\) 个形式级数 \(h_k(X_1, \ldots, X_n)\) 表示，其确定如下：对于 \(\alpha \in \mathbf{N}^n\) 和 \(\beta \in \mathbf{N}^m\)，令 \(c_{\alpha,\beta}\) 为形式级数 \(X^\alpha\) 中 \(\mathbf{f}^\beta = \prod f_j^{\beta_j}\) 的系数；则族 \((g_{k,\beta} c_{\alpha,\beta})_{\beta \in \mathbb{N}^m}\) 在 K 中可求和，其和就是 \(X^\alpha\) 中 \(h_k\) 的系数。

3.1.12. 假设 F 拟完备，且 \(\hat{\mathbf{E}}_{i}\) 是 \(\mathbf{E}_{i}\) 的完备化。在 \(\mathbf{E}_{1} \times \cdots \times \mathbf{E}_{n}\) 上取值于 F 的连续多项式可由连续性延拓为在 \(\hat{\mathbf{E}}_{1} \times \cdots \times \hat{\mathbf{E}}_{n}\) 上取值于 F 的多项式连续映射。由此得到从 \(j\) 到 \(\hat{\mathbf{P}}(\mathbf{E}_{1}, \ldots, \mathbf{E}_{n}; \mathbf{F})\) 的双射 \(\hat{\mathbf{P}}(\hat{\mathbf{E}}_{1}, \ldots, \hat{\mathbf{E}}_{n}; \mathbf{F})\)。若 \(f \in \mathcal{H}(\mathbf{E}_{1}, \ldots, \mathbf{E}_{n}; \mathbf{F})\)，则 \(j(f) \in \mathcal{H}(\hat{\mathbf{E}}_{1}, \ldots, \hat{\mathbf{E}}_{n}; \mathbf{F})\)，反之亦然。\(f\) 和 \(j(f)\) 的严格收敛指标相同。

3.1.13. 假设 K = R，但 F 赋有与其 R 向量空间结构相容的复向量空间结构。令 \(E_{i}^{C} = E_{i} \otimes_{R} C\)。若 \(y \in E_{i}^{C}\)，置：

\[
\| y \| = \inf \sum_ {k} | a _ {k} |. \| x _ {k} \|,
\]

下确界取遍满足 \((x_{k}, a_{k}) \in \mathrm{E}_{i} \times \mathbf{C}\) 的所有有限对族 \(y = \sum_{k} x_{k} \otimes a_{k}\)。于是得到复向量空间 \(E_{i}^{c}\) 上的一个范数；若同意将 x  \(E_{i}\) 与 \(x \in E_{i}\) 识别，则此范数延拓给定的 \(x \otimes 1\) 上的范数。设 \(h\) 是 \(E_{1} \times \cdots \times E_{n}\) 上取值于 F、关于多重次数 \(\alpha\) 齐次的 R-多项式连续映射；则存在唯一的在 \(\tilde{h}\) 上取值于 F 的 C-多项式连续映射 \(E_{1}^{c} \times \cdots \times E_{n}^{c}\)，延拓 h 且关于多重次数 \(\alpha\) 齐次。有

\[
\| \tilde {h} \| _ {\gamma} = \| h \| _ {\gamma}
\]

对复向量空间 F 上的每个连续半范数 γ。

\[
f = \sum_ {\alpha} f _ {\alpha} \in \mathcal {H} (\mathrm{E} _ {1}, \dots , \mathrm{E} _ {n}; \mathrm{F}), \text {   then   } \tilde {f} = \sum_ {\alpha} \tilde {f} _ {\alpha} \in \mathcal {H} (\mathrm{E} _ {1} ^ {\mathrm{c}}, \dots , \mathrm{E} _ {n} ^ {\mathrm{c}}; \mathrm{F}).
\]

级数 \(f\) 和 \(\tilde{f}\) 具有相同的严格收敛指标（且当 \(n = 1\) 时具有相同的严格收敛半径）。

反之，假设 K = C。令 \(E_{i}^{0}\) 和 \(F^{0}\) 为通过限制标量得到的 R 上的空间。若 \(f_{\alpha} \in \mathrm{P}_{\alpha}(\mathrm{E}_{1}, \ldots, \mathrm{E}_{n}; \mathrm{F})\)，则 \(f_{\alpha} \in \mathrm{P}_{\alpha}(\mathrm{E}_{1}^{0}, \ldots, \mathrm{E}_{n}^{0}; \mathrm{F}^{0})\)。若 \(f = \sum_{\alpha} f_{\alpha} \in \mathcal{H}(\mathrm{E}_{1}, \ldots, \mathrm{E}_{n}; \mathrm{F})\)，则形式级数 \(f^{0} = \sum_{\alpha} f_{\alpha} \in \hat{\mathrm{P}}(\mathrm{E}_{1}^{0}, \ldots, \mathrm{E}_{n}^{0}; \mathrm{F}^{0})\) 是收敛级数。f 和 \(f^{0}\) 的收敛指标（分别地，严格收敛指标）相同，并且对所有 \(f(x) = f^{0}(x)\)，有 \(x \in \tilde{\mathbf{C}}(f) = \tilde{\mathbf{C}}(f^{0})\)。

